% ============================================================
% §6.6 EQ5 — Physical Testbed End-to-End Validation
% IEEEtran double-column, no forced line breaks within paragraphs
% Label: sec:eq5 | Figure: \ref{fig:testbed-v6} | Table: \ref{tab:eq5-summary} / \ref{tab:testbed-hw}
% Five-sentence structure per experiment block
% ============================================================
\subsection{Physical Testbed Validation}
\label{sec:eq5}

Does \Name's dynamic priority adjustment translate from simulation to a
real RDMA fabric under contention?
We answer this question on a two-node single-bottleneck dumbbell testbed
(Fig.~\ref{fig:testbed-v6}) with a 6~Gbps iperf3 UDP background flow at
DSCP~P3 under strict-priority (SP) queueing, $c_i=1.2$ for tight jobs,
5-minute warmup, and no scheduler-side assignment logic (the closed-form
water-filling allocation of \S\ref{sec:eq1} degrades to a demand check:
unmet premium jobs receive strict priority, satisfied jobs maintain
best-effort SP handling).
We compare \Name (dynamic $P6{\rightarrow}P2$) against CRUX (static
P3) across four independent rounds (two orderings LL$\rightarrow$CX and
CX$\rightarrow$LL $\times$ two replications; a third replication, rep1,
was verified via md5sum to be a byte-identical copy and excluded).
Over the primary stable-contention window (epochs~7--11), \Name achieves
a mean slowdown of 1.11$\times$ against CRUX at 1.26--1.36$\times$
(Table~\ref{tab:eq5-summary}), a relative advantage of 13.7--16.8\%.
Three of the four independent rounds exhibit non-overlapping slowdown
intervals between the two schedulers; the remaining round (rep2\_r1)
produces overlapping intervals driven by an anomalous \Name elevation
that preserves the qualitative trend.  The overall replication rate is
thus 3/4 decisive non-overlap with 1/4 statistically indistinguishable.
The full window (epochs~7--14), which includes end-of-run contention
decay, serves as a sensitivity check: \Name remains at 1.09--1.21$\times$
while CRUX spans 1.10--1.34$\times$, with the CRUX low end
(1.10$\times$ at orig\_r1) an artifact of the competing job's completion.
\textbf{\Name's dynamic priority provides measurable end-to-end slowdown
reduction under realistic RDMA contention, consistent with the simulation
prediction of \S\ref{sec:eq2}.}

Two deployment insights from the testbed apply broadly.
First, the default DSCP-to-traffic-class mapping
($\texttt{trafficClass} = \mathtt{priority} \times 8$) silently inverts
software priority under SP scheduling\slash Table~\ref{tab:testbed-hw}(a):
the highest label (P6) maps to the lowest hardware class (tc:6) while the
lowest (P0) maps to tc:0.  The packet path remains fully functional,
making the misconfiguration invisible to standard connectivity checks;
only a full-class probe (DSCP~0--56) verified via NIC ${tx\_prio}$
counters (5.3~TB on prio3 vs.\ 0.049~TB on prio4 for P4)
revealed the inversion.  Secondly, the receiver NIC (node~226) does not
classify DSCP into separate priority queues\slash Table~\ref{tab:testbed-hw}(b):
all eight ${tx\_prio}$ counters record traffic exclusively on prio0
(9.79~TB), whereas the sender NIC (node~10.1) correctly distributes
across prio3,~prio4, and~prio6.  Consequently, per-priority enforcement
at this node sinks to the first switch hop.  These constraints are
properties of the specific hardware generation; they do not challenge
the scheduler design but must be accounted for in heterogeneous RDMA
deployments.

\begin{figure*}[htbp]
\centering
\includegraphics[width=\linewidth]{fig6_testbed.pdf}
\caption{Physical testbed comparison of \Name{} (dynamic, blue solid line)
and CRUX (static~P3, red dashed line) across four independent rounds.
Subfigure titles indicate run ordering (LongLiu$\rightarrow$CRUX or
CRUX$\rightarrow$LongLiu).  The gray band marks the stable-contention
window (epochs~7--11), the primary metric; epochs~7--14 are given as a
sensitivity reference.  rep2\_r1 (bottom-left) shows an anomalous
\Name{} elevation producing overlapping intervals (1/4 overlap).
Round~rep1 excluded via md5 byte-identity verification.
n=1 (single run per round; replications are independent experiment
repetitions, not seeds).}
\label{fig:testbed-v6}
\end{figure*}

\begin{table*}[htbp]
\centering
\caption{Stable contention window (epochs~7--11): per-round mean and
[minimum, maximum] slowdown for \Name{} and CRUX.  Rows marked with
$\dagger$ show non-overlapping intervals.  Rep1 excluded via md5
byte-identity verification.}
\label{tab:eq5-summary}
\footnotesize
\small
\begin{tabular}{lcccc}
\toprule
\textbf{Round} & \textbf{Order} & \textbf{Scheduler} & \textbf{Mean} & \textbf{[Min, Max]}\\
\midrule
\multirow{2}{*}{orig\_r1${}^\dagger$}
   & \multirow{2}{*}{LongLiu~$\rightarrow$~CRUX}
   & \Name  & 1.1075 & {[}1.0780, 1.1441{]} \\
   &                & CRUX  & 1.3022 & {[}1.2942, 1.3101{]} \\
\multirow{2}{*}{orig\_r2${}^\dagger$}
   & \multirow{2}{*}{CRUX~$\rightarrow$~LongLiu}
   & \Name  & 1.1069 & {[}1.0840, 1.1346{]} \\
   &                & CRUX  & 1.2824 & {[}1.2586, 1.3034{]} \\
\multirow{2}{*}{rep2\_r1}
   & \multirow{2}{*}{LongLiu~$\rightarrow$~CRUX}
   & \Name  & 1.2496 & {[}1.1378, 1.2909{]} \\
   &                & CRUX  & 1.2541 & {[}1.2478, 1.2618{]} \\
\multirow{2}{*}{rep2\_r2${}^\dagger$}
   & \multirow{2}{*}{CRUX~$\rightarrow$~LongLiu}
   & \Name  & 1.1123 & {[}1.0978, 1.1482{]} \\
   &                & CRUX  & 1.3365 & {[}1.2785, 1.3825{]} \\
\bottomrule
\end{tabular}
\end{table*}

% Evidence:
%   Stable window (epochs 7-11, primary):
%     FIGURE_REGISTRY/fig6_v6physical/fig6_data.csv, phase2_tight_stable rows
%     orig_r1: LL=1.1075 [1.0780..1.1441] vs CX=1.3022 [1.2942..1.3101]
%     orig_r2: LL=1.1069 [1.0840..1.1346] vs CX=1.2824 [1.2586..1.3034]
%     rep2_r1: LL=1.2496 [1.1378..1.2909] vs CX=1.2541 [1.2478..1.2618] (overlap)
%     rep2_r2: LL=1.1123 [1.0978..1.1482] vs CX=1.3365 [1.2785..1.3825]
%     Advantage range: 13.7% (orig_r2) -- 16.8% (rep2_r2)
%   Full window (epochs 7-14, sensitivity):
%     same CSV, phase2_tight_full rows; LL 1.09--1.21, CX 1.10--1.34
%   Deprecated: rep1 md5-identical to orig (9438c116 / 3f8170eb), excluded
%   Deployment evidence:
%     PAPER_EVIDENCE/00_documents/t3_data_table.csv (DSCP probe / 226 probe / tc_map)
%     Table~\ref{tab:testbed-hw}=FIGURE_REGISTRY/t3_topo/tab_t3.tex
